\documentclass{article}
\usepackage[a4paper, total={6.5in, 10in}]{geometry}
\setlength{\parindent}{0pt}
\usepackage{tcolorbox}
\tcbset{colback=white!94!black, colframe=black!60!white, coltitle = white, boxrule=0.8pt, arc=2mm}
\newtcolorbox{boxs}[2][]{title= #2,#1}
\usepackage{datetime2}
\usepackage[british]{babel}
\usepackage{amsfonts}
\usepackage{amsmath}

\title{\textbf{Advanced Calculus HW 7}}
\author{Robert O'Connell}
\date{\today}
\begin{document}
\maketitle
\vspace{5mm}
\subsection*{Question 1}
We know that the tangent plane to \(F\) at \((x_0, y_0,z_0)\) is given by
\[
F(x_0,y_0,z_0) + F_x(x_0,y_0)(x-x_0) + F_y(x_0,y_0)(y-y_0)  + F_z(z_0,y_0,z_0)(z-z_0) = 0
\]

And its normal line is given by \(\vec{l}(t) = (x_0,y_0,z_0) + t(\nabla f(x_0,y_0,z_0))\)

\subsubsection*{(a)}
\[
F(x,y,z) = z - 2 - \frac{1}{2}\ln(x^2 + y^2)
\]
\[
F(-1,0,2) = 0
\]
\[
\nabla F(x,y,z) = (\frac{-x}{x^2 + y^2}, \frac{-y}{x^2 + y^2},  1)
\]
\[
\nabla F(-1,0,2) = (1,0,1)
\]

Then the tangent plane is given by
\[
(x+1) + (z-2) = 0
\]
\[
x+z-1 =0
\]

And the normal line is 
\[
\vec{l}(t) = (-1,0,2) + t(1,0,1)
\]

\subsubsection*{(b)}
\[
F(x,y,z) = z - \frac{1}{2}x^7y^{-2} -1
\]
\[
F(2,4,3) = 0
\]
\[
\nabla F(x,y,z) = (-\frac{7}{2}x^6y^{-2}, x^7y^{-3},  1)
\]
\[
\nabla F(2,4,3) = (-14,2,1)
\]

Then the tangent plane is given by
\[
-14(x-2) + 2(y-4)+ z - 3 = 0
\]
\[
-14x + 2y + z + 17 =0
\]

And the normal line is 
\[
\vec{l}(t) = (2,4,3) + t(-14,2,1)
\]





\subsection*{Question 2}
\subsubsection*{(a)}
\[
w = f(\rho) = f(\sqrt{(x^2 + y^2 + z^2)})
\]
\[
\frac{\partial w}{\partial x} = \frac{dw}{d\rho} \frac{\partial \rho}{\partial x} 
\]

with \(\frac{\partial \rho}{\partial x} = \frac{x}{\sqrt{x^2 + y^2 + z^2}} = \frac{x}{\rho}\), since \(\rho\) is similar
across \(x,y\) and \(z\), their partial derivatives are similar

Then
\begin{align*}
    \left(\frac{\partial w}{\partial x}\right)^2 + \left(\frac{\partial w}{\partial y}\right)^2 + \left(\frac{\partial w}{\partial z}\right)^2 &= 
    \left(\frac{df}{d \rho} \frac{x}{\rho}\right)^2 + \left(\frac{df}{d \rho} \frac{y}{\rho}\right)^2 + \left(\frac{df}{d \rho} \frac{z}{\rho}\right)^2 \\
    &= \left(\frac{df}{d \rho}\right)^2\left(\frac{x^2}{\rho^2} + \frac{y^2}{\rho^2} + \frac{z^2}{\rho^2}\right) \\
    &= \left(\frac{df}{d \rho}\right)^2\left(\frac{1}{\rho^2}\right)(x^2 + y^2 + z^2) \\
    &= \left(\frac{df}{d \rho}\right)^2\left(\frac{\rho^2}{\rho^2}\right) \\
    &= \left(\frac{df}{d \rho}\right)^2
\end{align*}

Also
\[
\frac{dw}{d\rho} = \frac{df}{d \rho}
\]

Then 
\[
\left(\frac{\partial w}{\partial x}\right)^2 + \left(\frac{\partial w}{\partial y}\right)^2 + \left(\frac{\partial w}{\partial z}\right)^2 = \left(\frac{dw}{d \rho}\right)^2
\]





\subsubsection*{(b)}
\[
w = f(x,y,g(x,y))
\]
\[
\frac{\partial w}{\partial x} = \frac{\partial f}{\partial x} + \frac{\partial f}{\partial g}\frac{\partial g}{\partial x}
\]
\[
\frac{\partial w}{\partial y} = \frac{\partial f}{\partial y} + \frac{\partial f}{\partial g}\frac{\partial g}{\partial y}
\]


\subsubsection*{(c)}

\[
w = f(x,y,z)
\]
\[
\frac{dw}{dt} = \frac{\partial f}{\partial x} \frac{dx}{dt} + \frac{\partial f}{\partial y} \frac{dy}{dt} + \frac{\partial f}{\partial z} \frac{dz}{dt}
\]
\[
= \left(\frac{6x}{3x^2 - 2y + 4z^3}\right)\left(\frac{1}{2\sqrt{t}}\right) + \left(-\frac{2}{3x^2 - 2y + 4z^3}\right)
\left(\frac{2}{3}t^{\frac{-1}{3}}\right) + \left(\frac{12z^2}{3x^2 - 2y + 4z^3}\right)\left(-\frac{1}{t^2}\right)
\]
\[
= \left(-\frac{2}{3x^2 - 2y + 4z^3}\right)\left(-\frac{3x}{2\sqrt{t}} + \frac{2}{3t^{\frac{1}{3}}} + \frac{6z^2}{t^2}\right)
\]


\subsection*{Question 3}
\subsubsection*{(a)}
First normalise \(\vec{u}\)
\[
\hat{u} = \left(\frac{1}{\sqrt{14}}, \frac{2}{\sqrt{14}}, -\frac{3}{\sqrt{14}}\right)
\]

\[
\nabla f(x,y,z) = \left(-\frac{y}{(x^2+z^2)^2},\frac{1}{x+z} , -\frac{y}{(x^2+z^2)^2}\right)
\]

Then the directional derivative is given by \(\hat{u} \cdot \nabla f(x_0,y_0,z_0)\)
\[
\left(-\frac{y_0}{\sqrt{14}(x_0^2+z_0^2)^2},\frac{2}{\sqrt{14}(x_0+z_0)} , \frac{3y_0}{\sqrt{14}(x_0^2+z_0^2)^2}\right)
\]

\subsubsection*{(b)}
\[
f(x,y,z) = y^2z \tan^3x
\]
\[
f_x(x,y,z) = 3y^2z \tan^2x \sec^2x
\]
\[
f_y(x,y,z) = 2yz \tan^3x
\]
\[
f_z(x,y,z) = y^2 \tan^3x
\]

\[
\Rightarrow \nabla f(x,y,z) = (3y^2z \tan^2x \sec^2x, 2yz \tan^3, y^2 \tan^3)
\]
\[
\Rightarrow \nabla f(\pi / 4, -2, 1) = (24,-4,4)
\]
\end{document}